\documentclass[aspectratio=43]{beamer}
\usepackage{luatexja}
\usepackage{hyperref,listings,tabularx,tikz,pifont}
\usetikzlibrary{arrows.meta,positioning,calc,fit}
\tikzset{str/.style={rounded corners, minimum height=1cm,align=center}}
\usepackage[symbol]{footmisc}
\newcommand{\cmark}{{\Large\bfseries\color{green!50!black}\ding{51}}}
\newcommand{\xmark}{{\Large\bfseries\color{red!50!black}\ding{55}}}
\hypersetup{
    colorlinks=true,
    linkcolor=black,
    urlcolor=blue
}
\usetheme{Montpellier}
\usecolortheme[RGB={160,82,45}]{structure}
\setbeamertemplate{navigation symbols}{}
\renewcommand{\familydefault}{\rmdefault}
\renewcommand{\figurename}{図\thefigure}
\renewcommand{\tablename}{表\thetable}
\renewcommand{\lstlistingname}{src.}
\renewcommand{\thefootnote}{\fnsymbol{footnote}}
\renewcommand{\thempfootnote}{\fnsymbol{mpfootnote}}
\bibliographystyle{junsrt}
\setbeamertemplate{bibliography item}[text]
\setbeamertemplate{section in toc}{
    {\color{structure}\inserttocsectionnumber.}\ \inserttocsection\vspace{.3em}
}

\setbeamertemplate{subsection in toc}{
    \hspace{1.2em}{\color{structure}\rule[.3ex]{3pt}{3pt}}\ \inserttocsubsection\par\vspace{.3em}
}
\setbeamertemplate{footline}{
    \leavevmode
    \hbox{
        \begin{beamercolorbox}[wd=.95\paperwidth,ht=1ex,right,dp=1.5ex,rightskip=0em]{bg=gray}
            \insertframenumber{} / \inserttotalframenumber
        \end{beamercolorbox}
    }
}
\setbeamercolor{block title}{use=structure,fg=white,bg=structure.fg!75!black}
\setbeamercolor{block title example}{parent=example text, fg=white,bg=example text.fg!75!black}
\setbeamercolor{block title alerted}{parent=alerted text, fg=white,bg=alerted text.fg!75!black}
\setbeamercolor{block body}{parent=normal text,use=block title,bg=block title.bg!10}
\setbeamercolor{block body alerted}{parent=normal text,use=block title alerted,bg=block title alerted.bg!10}
\setbeamercolor{block body example}{parent=normal text,use=block title example,bg=block title example.bg!10}
\setbeamertemplate{enumerate item}[default]
\setbeamertemplate{itemize item}[triange]
\newcolumntype{C}{>{\centering\arraybackslash}X}
\newcolumntype{R}{>{\raggedright\arraybackslash}X}
\newcolumntype{L}{>{\raggedleft\arraybackslash}X}
\lstdefinelanguage{docker-compose-2}{
    keywords={version, volumes, services},
    morekeywords={FROM,RUN,CMD,EXPOSE,ENV},
    keywords=[2]{image, environment, ports, container_name, ports, links, build},
    sensitive=false,
    comment=[l]{\#},
    morestring=[b]',
    morestring=[b]"
}
\lstset{
    %枠外に行った時の自動改行
    breaklines = true,
    %自動改行後のインデント量(デフォルトでは20[pt])
    breakindent = 10pt,
    %標準の書体
    basicstyle = \ttfamily\small,
    %コメントの書体
    commentstyle = {\scriptsize\ttfamily \color[cmyk]{1,0.4,1,0}},
    %関数名等の色の設定
    classoffset = 0,
    %キーワード(int, ifなど)の書体
    keywordstyle = {\bfseries\ttfamily \color[cmyk]{0,1,0,0}},
    %表示する文字の書体
    stringstyle = {\ttfamily \color[rgb]{0,0,1}},
    %枠 tは上に線を記載, Tは上に二重線を記載
    %他オプション：leftline，topline，bottomline，lines，single，shadowbox
    frame = single,
    %frameまでの間隔(行番号とプログラムの間)
    framesep = 5pt,
    %行番号の位置
    numbers = left,
    %行番号の間隔
    stepnumber = 1,
    %行番号の書体
    numberstyle = \scriptsize,
    %タブの大きさ
    tabsize = 4,
    %キャプションの場所(tbならば上下両方に記載)
    captionpos = t,
    % 左余白
    xleftmargin=1em,
}
\makeatletter
\@addtoreset{footnote}{page}
\makeatother
\newcommand{\ftitle}{\thesection.\thesubsection\ \subsecname}
\newcommand{\tocc}{
    \begin{frame}[t]{\thesection.\ \secname}
        \begin{columns}[t]
            \begin{column}{.48\textwidth}
                \tableofcontents[sections={1-3},currentsection,sectionstyle=show/shaded,subsectionstyle=show/show/shaded]
            \end{column}
            \begin{column}{.48\textwidth}
                \tableofcontents[sections={4},currentsection,sectionstyle=show/shaded,subsectionstyle=show/show/shaded]
            \end{column}
        \end{columns}
    \end{frame}
}
\title{仮想化技術とクラウドコンピューティング}
\subtitle{AWS第1回}
\author{吉川 恭平\and 塩澤 康志\and 福島 康記\and 溝口 洸熙}
\institute{Group C5}
\date{発表日}
\begin{document}
\begin{frame}
    \maketitle
\end{frame}
\begin{frame}[t]{お品書き}
    \begin{columns}[t]
        \begin{column}{.48\textwidth}
            \tableofcontents[sections={1-3}]
        \end{column}
        \begin{column}{.48\textwidth}
            \tableofcontents[sections={4}]
        \end{column}
    \end{columns}
\end{frame}
\section{仮想化技術と仮想マシン}
\tocc
\newcommand{\ie}{\tikz{\draw[arrows={-Latex[line width=1pt,fill=white,length=10pt]}](0,0)--(1,0)}\ }
\begin{frame}[t]{\thesection.\ \secname}
    \begin{block}{仮想化}
        ``コンピューターの物理的資源を論理的に分割して，それぞれ独立並列した状態で利用できるようにすること．
        1台のサーバーで，複数の基本ソフトを独立並列に動作させるサーバー仮想化など．''\hfill\cite{スーパー大辞林}
    \end{block}
    \ie 1台の計算機で複数のOSやアプリケーションなどを並列に動作させること．
    \vspace{.5em}
    \begin{block}{Virtual Machine（仮想機械，仮想マシン）}
        ``あるコンピューターシステムの動作を，別システムで再現するソフトウエア．また，そのような動作環境．あるOSの動作を別のOS上で再現する場合など．バーチャルマシン．VM．''\hfill\cite{スーパー大辞林}
    \end{block}
    \ie 仮想化するためのソフトウェアや動作環境．
\end{frame}
\section{仮想化の方式}
\tocc
\subsection{ホスト型}
\begin{frame}[t]{\ftitle}
    ハードウェアの中のOS上に，土台となる仮想ソフトウェアをインストールし，仮想化ソフトウェアで仮想マシンを稼働させる．
    \begin{figure}[b]
        \centering
        \begin{tikzpicture}
            \node[str,fill=gray!60,text width=.9\textwidth](hw){ハードウェア};
            \node[str,fill=gray!50,text width=.9\textwidth,above=.1cm of hw](hos){ホストOS};
            \node[str,fill=gray!40,text width=.9\textwidth,above=.1cm of hos](hsw){ホスト型仮想化ソフトウェア};
            \node[str,fill=gray!30,text width=.43\textwidth,above=.1cm of hsw.north west,anchor=south west](gos1){ゲストOS};
            \node[str,fill=gray!30,text width=.43\textwidth,above=.1cm of hsw.north east,anchor=south east](gos2){ゲストOS};
            \node[str,fill=gray!20,text width=.2\textwidth,above=.1cm of gos1.north west,anchor=south west](sw1){\tiny アプリケーション};
            \node[str,fill=gray!20,text width=.2\textwidth,above=.1cm of gos1.north east,anchor=south east](sw2){\tiny アプリケーション};
            \node[str,fill=gray!20,text width=.2\textwidth,above=.1cm of gos2.north west,anchor=south west](sw3){\tiny アプリケーション};
            \node[str,fill=gray!20,text width=.2\textwidth,above=.1cm of gos2.north east,anchor=south east](sw4){\tiny アプリケーション};
        \end{tikzpicture}
    \end{figure}
    % \hyperlink{ハイパーバイザ型}{\beamergotobutton{ハイパーバイザ型}}\hypertarget{ホスト型}{}
\end{frame}
\begin{frame}[t]{\ftitle}
    \begin{exampleblock}{ホスト型仮想化ソフトウェア 例}
        \begin{minipage}[b]{.8\textwidth}
            \begin{itemize}
                \setlength{\itemsep}{1em}
                \item VMware Workstation Player
                \item VMware Fusion
                \item Oracle VM Virtualbox
            \end{itemize}
        \end{minipage}
        \begin{minipage}[b]{.15\textwidth}
            \centering
            \includegraphics[keepaspectratio,width=\textwidth]{virtualbox_logo.png}\\
            {\tiny Virtualbox\cite{VMBox}}
        \end{minipage}
    \end{exampleblock}
    \begin{minipage}[t]{.48\textwidth}
        \textbf{メリット}
        \begin{itemize}
            \item 既存マシンが利用できる点．
            \item 仮想化に必要なソフトウェアが扱いやすい．
        \end{itemize}
    \end{minipage}
    \begin{minipage}[t]{.48\textwidth}
        \textbf{デメリット}
        \begin{itemize}
            \item ホストOSを動作させるための物理リソースが必要．
        \end{itemize}
    \end{minipage}\\
    \hfill\cite{itmanage}
\end{frame}
\subsection{ハイパーバイザ型}
\begin{frame}[t]{\ftitle}
    ハイパーバイザとは「仮想化のためのOS」のようなもの．
    \begin{figure}[b]
        \centering
        \begin{tikzpicture}
            \node[str,fill=gray!60,text width=.9\textwidth](hw){ハードウェア};
            \node[str,fill=gray!45,text width=.9\textwidth,above=.1cm of hw,minimum height=2.1cm](hpv){ハイパーバイザ};
            \node[str,fill=gray!30,text width=.43\textwidth,above=.1cm of hpv.north west,anchor=south west](gos1){ゲストOS};
            \node[str,fill=gray!30,text width=.43\textwidth,above=.1cm of hpv.north east,anchor=south east](gos2){ゲストOS};
            \node[str,fill=gray!20,text width=.2\textwidth,above=.1cm of gos1.north west,anchor=south west](sw1){\tiny アプリケーション};
            \node[str,fill=gray!20,text width=.2\textwidth,above=.1cm of gos1.north east,anchor=south east](sw2){\tiny アプリケーション};
            \node[str,fill=gray!20,text width=.2\textwidth,above=.1cm of gos2.north west,anchor=south west](sw3){\tiny アプリケーション};
            \node[str,fill=gray!20,text width=.2\textwidth,above=.1cm of gos2.north east,anchor=south east](sw4){\tiny アプリケーション};
        \end{tikzpicture}
    \end{figure}
    % \hyperlink{ホスト型}{\beamergotobutton{ホスト型}}\hypertarget{ハイパーバイザ型}{}
\end{frame}
\begin{frame}[t]{\ftitle}
    \begin{exampleblock}{ハイパーバイザ 例}
        \begin{itemize}
            \item VMware ESXi
            \item Linux KVM
            \item Microsoft Hyper-V
        \end{itemize}
    \end{exampleblock}
    \begin{minipage}[t]{.48\textwidth}
        \textbf{メリット}
        \begin{itemize}
            \item システム全体の観点から見てリソースの使用効率がよい．
            \item ホストOSが不要でハードウェアを直接制御が可能．
        \end{itemize}
    \end{minipage}
    \begin{minipage}[t]{.48\textwidth}
        \textbf{デメリット}
        \begin{itemize}
            \item 仮想化環境の高度な管理を実現するツールが標準装備されていない場合がある．
        \end{itemize}
    \end{minipage}\\
    \hfill\cite{itmanage}
\end{frame}
\subsection{コンテナ}
\begin{frame}[t]{\ftitle}
    ``アプリケーションを実行するための領域（ユーザ空間）を複数に分割して利用するもの''\cite{itmanage}．
    \begin{figure}[b]
        \centering
        \begin{tikzpicture}
            \node[str,fill=gray!60,text width=.9\textwidth](hw){ハードウェア};
            \node[str,fill=gray!50,text width=.9\textwidth,above=.1cm of hw](hos){ホストOS};
            \node[str,fill=gray!40,text width=.9\textwidth,above=.1cm of hos](hsw){Docker または Kubernetes};
            \node[str,fill=gray!30,text width=.2\textwidth,above=.1cm of hsw.north west,anchor=south west](gos1){ミドルウェア};
            \node[str,fill=gray!30,text width=.2\textwidth,above=.1cm of hsw.north east,anchor=south east](gos2){ミドルウェア};
            \node[str,fill=gray!20,text width=.2\textwidth,above=.1cm of gos1.north west,anchor=south west](sw1){\tiny アプリケーション};
            \node[str,fill=gray!20,text width=.2\textwidth,above=.1cm of gos2.north east,anchor=south east](sw4){\tiny アプリケーション};
            \node[inner sep=.5mm,fit={(gos1)(sw1)},draw,thick,dotted,rounded corners](wrap1){};
            \node[inner sep=.5mm,fit={(gos2)(sw4)},draw,thick,dotted,rounded corners](wrap2){};
            \node at ($(wrap1)!0.5!(wrap2)+(0,.5cm)$)(cont){コンテナ};
            \draw[-latex](cont.west)--(wrap1.east);
            \draw[-latex](cont.east)--(wrap2.west);
        \end{tikzpicture}
    \end{figure}
\end{frame}
\section{Docker（Option）}\label{chap:docker}
\tocc
\subsection{コンテナ立ち上げと利用}
\begin{frame}[t]{\ftitle}
    \begin{itemize}
        \setlength{\itemsep}{1em}
        \item コンテナを利用する．
              \begin{enumerate}
                  \setlength{\itemsep}{.5em}
                  \item イメージを作成（\texttt{build}）または，取得（\texttt{docker pull}）する．\\
                        イメージの作成には\texttt{Dockerfile} を書く．
                  \item コンテナを作成して起動（\texttt{run}）する．
                  \item コンテナ内でコマンドを実行（\texttt{exec}）する．
              \end{enumerate}
        \item コンテナ利用後．
              \begin{enumerate}
                  \setlength{\itemsep}{.5em}
                  \item コンテナを停止（\texttt{stop}）する．
                  \item コンテナを削除（\texttt{rm}）する．
                  \item イメージを削除（\texttt{rmi}）する．
              \end{enumerate}
    \end{itemize}
\end{frame}
\subsection{イメージ作成・取得}
\begin{frame}[containsverbatim,t]{\ftitle}
    \begin{itemize}
        \item \texttt{Dockerfile}
              \begin{lstlisting}[language=docker-compose-2]
FROM mysql:latest
            \end{lstlisting}
              \begin{lstlisting}[backgroundcolor={\color{black}},basicstyle={\color{white}\ttfamily},]
$ docker build ./ -t mysql
                            \end{lstlisting}
        \item レジストリからイメージを取得（Pull）する．
              \begin{lstlisting}[backgroundcolor={\color{black}},basicstyle={\color{white}\ttfamily},]
$ docker pull mysql:latest
            \end{lstlisting}
    \end{itemize}
\end{frame}
\subsection{実演}
\begin{frame}{\ftitle}
    \includegraphics[keepaspectratio,width=\textwidth]{docker_logo.png}\\\vspace{1em}
    \hfill\cite{docker}
\end{frame}
\section{クラウドコンピューティング}\hypertarget{sec:クラウドコンピューティング}{}
\tocc
\subsection{クラウドとは}
\begin{frame}[t]{\ftitle}
    \begin{block}{クラウドコンピューティング}
        クラウドコンピューティング（以下，クラウド）とは，コンピュータリソース尾の利用形態．
        コンピュータの計算リソースやストレージ領域，アプリケーションによる処理をネットワーク経由で提供する．\hfill\cite{2015amazon}
    \end{block}
    クラウドは以下の4つに分類されることが多い．
    \begin{itemize}
        \setlength{\itemsep}{1em}
        \item IaaS（Infrastructure as a Service）
        \item PaaS（Platform as a Service）
        \item SaaS（Software as a Service）
        \item DaaS（Desktop as a Service）% VDIとの違い
    \end{itemize}
\end{frame}
\subsection{クラウドの分類}
\begin{frame}[t]{\ftitle}
    \begin{itemize}
        \item \textbf{IaaS}\\
              仮想サーバやストレージ，ネットワークサービスなどをインターネット経由で提供する．\\（AWS，Azure，GCP\footnote{Google Cloud Pratform}，IBM Cloudなど）
        \item \textbf{PasS}\\
              ミドルウェアをサービスとして提供する．
              OSとミドルウェアの管理はサービス提供者で行われ，ユーザはミドルウェアを直接操作できる．\\（AWS，Azure，GCPなど）
        \item \textbf{SaaS}\\
              ソフトウェアをサービスとして提供する．\\（Slack，Zoom，Google Calendarなど）
    \end{itemize}
    \hfill\cite{2015amazon}
\end{frame}
\begin{frame}[t]{\ftitle}
    \begin{itemize}
        \item \textbf{DaaS}\\
              ``クラウドサービスとして提供される仮想デスクトップ環境を組織として利用するしかけ．''\hfill\cite{ISディジタル辞典DaaS}\\
              %   \begin{itemize}
              %       \item[{\bfseries\Large\cmark}] \underline{クライアントを安価に抑えることができる．}（クライアントが依頼した処理をサーバが処理するため）
              %       \item[{\bfseries\Large\xmark}] \underline{ネットワークネットワーク負荷が大きくなる．}（ディスクトップ環境を提供するため）
              %   \end{itemize}
              （Azure Virtual Desktopなど）\\
              \hrulefill
        \item \textbf{VDI}（Virtual Desktop Infrastructure）\\
              ``クライアントPCのディスクトップサーバ上に仮想化して集約し，稼働させる仕組み．''\hfill\cite{ISディジタル辞典VDI}
              \item[\textbf{\large !}]DaaSはVDIの一種である．
    \end{itemize}
\end{frame}
\subsection{代表的なIaaS}
\begin{frame}[t]{\ftitle}
    \begin{itemize}
        \setlength{\itemsep}{1em}
        \item \underline{AWS}
        \item Microsoft Azure
        \item Google Cloud Pratform
        \item さくらのクラウド（さくらインターネット）
    \end{itemize}
\end{frame}
\begin{frame}[allowframebreaks]{Reference}
    \bibliography{bib}
\end{frame}
\end{document}